% papers/partitura.tex
% Partitura do Tradutor --- desenhada pelo tradutor .tex↔PDF (tex.wasm).
% Alinhada ao gabarito LilyPond (partitura/render/tradutor-completo.pdf):
%   StaffGroup, 8 compassos × 4 tempos, mesmo relógio em todas as vozes.
% Léxico: \note → ♪ (U+266A). LilyPond em partitura/ só grava.
%
%   tests/tex papers/partitura.tex saida.pdf
\documentclass[11pt,a4paper]{article}
\usepackage[utf8]{inputenc}
\usepackage[T1]{fontenc}
\usepackage[portuguese]{babel}
\usepackage{amsmath,amssymb,amsthm}
\usepackage[margin=2.5cm]{geometry}
\usepackage{booktabs}
\usepackage{xcolor}
\usepackage[hidelinks]{hyperref}
\newtheorem{definicao}{Definição}
\newtheorem{observacao}{Observação}
\newcommand{\code}[1]{\texttt{#1}}

% ─── glifos do relógio (mesmo padrão do gabarito emite.c / LilyPond) ─────────
\newcommand{\N}{\note}                 % ataque
\newcommand{\Rr}{\ensuremath{\cdot}}   % silêncio / descanso
% um compasso = 4 tempos (células da tabela alinhadas entre vozes)
\newcommand{\barA}{\N\N\N\N}           % quatro ataques (metr./cordas/...)
\newcommand{\barB}{\N\Rr\N\Rr}         % ataque-pausa (batuta/percussão)
\newcommand{\barC}{\Rr\N\Rr\N}         % pausa-ataque (batuta alternada)
\newcommand{\barM}{\N\Rr\Rr\Rr}        % semibreve ≈ ataque + 3 pausas (maestro)

% ─────────────────────────────────────────────────────────────────────────────
%  papers/gkcapa.tex — a capa do Reino, mínima, para os papers em `article`.
%
%  O estilo.tex completo é do `report` (títulos de capítulo, skak, …). Aqui fica
%  só o que a capa precisa, para o paper continuar a compilar sozinho com
%  pdflatex. O tradutor próprio (tests/tex) carrega o estilo.tex da raiz / o
%  symlink papers/estilo.tex e avalia o mesmo `\gkcapa`.
% ─────────────────────────────────────────────────────────────────────────────
\definecolor{ouro}{HTML}{B8912F}
\definecolor{ouroclaro}{HTML}{F3E9CE}
\definecolor{tinta}{HTML}{1E1B16}
\definecolor{regua}{HTML}{6E6858}
\providecommand{\gknota}{\fontsize{7.62}{11.05}\selectfont}
\providecommand{\gktexto}{\fontsize{10.50}{15.22}\selectfont}
\providecommand{\gksub}{\fontsize{12.33}{17.87}\selectfont}
\providecommand{\gksec}{\fontsize{14.47}{20.98}\selectfont}
\providecommand{\gkcap}{\fontsize{16.99}{24.63}\selectfont}
\providecommand{\gktit}{\fontsize{23.42}{33.95}\selectfont}
\providecommand{\Prj}{\mathbb{P}}
\providecommand{\Trf}{\mathcal{T}}
\providecommand{\gkcapa}[3]{%
\title{\vspace{-0.6cm}
{\color{ouro}\rule{\textwidth}{1.4pt}}\\[6mm]
\textbf{\gktit\color{tinta} Reino Dourado}\\[3mm]{\gksec\itshape\color{regua} Golden Kingdom}\\[8mm]
{\gkcap\scshape\color{ouro} #1}\\[5mm]
{\gksub\color{regua} #2}\\[2mm]
{\gktexto\color{regua}\itshape #3}\\[7mm]
{\color{ouro}\rule{\textwidth}{0.6pt}}\\[9mm]{\gknota\color{regua}\begin{minipage}{\textwidth}\setlength{\parindent}{0pt}%
\textbf{Aviso de Isenção Algorítmica:} O universo, as leis e as entidades contidas nestas páginas ---
incluindo felinos matriciais, aves de rapina algébricas e amuletos de controle de fluxo --- são
construções de pura ficção formal. Esta obra não descreve a realidade física, não serve como
engenharia reversa do cosmos e não constitui doutrina religiosa. Qualquer semelhança com bugs reais do seu
sistema operacional é mera coincidência (ou falta de reza). Prossiga por sua própria conta e
risco.\end{minipage}}}
\author{}\date{}}

\begin{document}
\gkcapa{Partitura do Tradutor}
       {$\Pi_{\mathrm{tradutor}}\longmapsto\Pi_{\mathrm{musical}}$}
       {StaffGroup no tradutor --- alinhado ao gabarito LilyPond}
\maketitle

\begin{abstract}
\noindent Especificação musical do sistema, composta pelo tradutor
\texttt{.tex}$\leftrightarrow$PDF. O léxico mapeia \verb|\note| ao glifo ♪;
o \textbf{StaffGroup} é uma tabela cujas colunas são o mesmo tick em todas
as vozes --- a mesma geometria do gabarito
\code{partitura/render/tradutor-completo.pdf}. Objectivo:
\[
\boxed{
\Pi_{\mathrm{tradutor}}
\;\longmapsto\;
\Pi_{\mathrm{musical}}
\quad(\text{todas as vozes sincronizadas}).
}
\]
LilyPond (\code{partitura/emite.c}) só grava; não inventa campos.
\end{abstract}

\section{Objectivo}

\[
\boxed{
\text{Tradutor}
\;\xrightarrow{\Pi}\;
\text{Partitura completa}
\;\xrightarrow{\text{tick}}\;
\{G_0,\ldots,G_n\}
\;\xrightarrow{\text{Maestro}}\;
\text{Orquestra}
}
\]
Estado $\to$ $\Pi$ $\to$ vozes paralelas $\to$ execução sincronizada.

\section{Gabarito e relógio}

\begin{definicao}[StaffGroup Tiffany]
Oito vozes $\times$ oito compassos $\times$ quatro tempos. Cada
\textbf{coluna} da tabela é um compasso; a mesma coluna em todas as
linhas é o \textbf{mesmo instante} (como as barras do StaffGroup LilyPond).
\end{definicao}

\begin{observacao}
Ritmos = os de \code{emite tradutor}: Metrónomo e naipes em semínimas
contínuas; Batuta e Percussão em ataque-pausa; Maestro em semibreves
(um ataque por compasso); Lei~8 no ciclo de oito graus.
\end{observacao}

\section{A partitura do sistema}

\noindent{\small Dim$=$8 · Alcance$=$3 · Lado$=$0 · Iface$=$6 · bpm$=$72 ·
8 compassos · \texttt{\string\time} 4/4}\par\medskip

\noindent
\begin{tabular}{@{}lcccccccc@{}}
voz & 1 & 2 & 3 & 4 & 5 & 6 & 7 & 8 \\
\midrule
Metronomo
 & \barA & \barA & \barA & \barA
 & \barA & \barA & \barA & \barA \\
Batuta
 & \barB & \barC & \barB & \barC
 & \barB & \barC & \barB & \barC \\
Maestro
 & \barM & \barM & \barM & \barM
 & \barM & \barM & \barM & \barM \\
Cordas (1)
 & \barA & \barA & \barA & \barA
 & \barA & \barA & \barA & \barA \\
Madeiras (2)
 & \barA & \barA & \barA & \barA
 & \barA & \barA & \barA & \barA \\
Metais (4)
 & \barA & \barA & \barA & \barA
 & \barA & \barA & \barA & \barA \\
Percussao (8)
 & \barB & \barB & \barB & \barB
 & \barB & \barB & \barB & \barB \\
Lei 8
 & \barA & \barA & \barA & \barA
 & \barA & \barA & \barA & \barA \\
\end{tabular}

\begin{observacao}
Confrontar com \code{partitura/partitura.pdf} (gravação LilyPond): as
\textbf{colunas} batem com as barras do sistema; a alteração estrutural
aparece no mesmo instante em todas as vozes. O pentagrama tipográfico
fica na gravação; aqui o tradutor realiza o sincronismo.
\end{observacao}

\section{Pipeline}

\[
\text{Corpo}
\;\to\;
\Pi
\;\to\;
\text{tradutor (StaffGroup + \note)}
\;\parallel\;
\text{emite.c}\to\texttt{.ly}\to\text{LilyPond}.
\]

\section{Fecho}

\begin{observacao}
Uma única partitura, todas as vozes, um único relógio, uma única
regência. LilyPond grava; a partitura é estrutura; o tradutor é a fonte.
\end{observacao}

\end{document}
